\section{El RL Environment es arquitectura de producto: tareas, herramientas, recompensas y recuperación definen juntas el comportamiento}

La tercera parte de la clase condensa su título en una sola frase: \textbf{how you construct RL environment and rewards is how your product will work}. No es retórica. Si el entorno de entrenamiento solo admite respuestas de texto, el modelo no aprenderá de la nada a usar herramientas de forma confiable; si la reward solo considera la cadena final, el modelo puede ignorar los permisos, el costo y los efectos secundarios irreversibles del proceso; si el entorno no contempla la recuperación ante fallos, la política no tiene oportunidad de aprender cuándo detenerse, pedir ayuda o revertir.

\lecturefigure{slide-50-rl-environment-product.jpg}{Tesis central: la construcción del RL environment y de la reward determina el comportamiento del producto}{00:33:15--00:35:51}

\begin{importantbox}{Primer uso: MDP, POMDP y contract del entorno}
El Markov Decision Process (MDP) estándar se escribe como
\[
\mathcal{M}=(\mathcal{S},\mathcal{A},P,R,\gamma),
\]
donde $\mathcal{S}$ es el estado, $\mathcal{A}$ la acción, $P$ la transición de estados, $R$ la recompensa y $\gamma$ el descuento. Un producto real se parece más a un Partially Observable MDP (POMDP): el modelo solo ve la observation $o_t$; la intención del usuario, el estado de las herramientas y las consecuencias futuras no son completamente visibles. El contract del entorno también debe declarar permisos, presupuesto, terminación, registros, confirmación humana y recuperación.
\end{importantbox}

El retorno de la trayectoria es
\[
G_t=\sum_{k=0}^{T-t}\gamma^k r_{t+k}.
\]
En las tareas largas, lo decisivo no es la fórmula, sino de dónde proviene $r_t$. Las pruebas de código, el estado de la base de datos, la confirmación del usuario y el judge humano tienen distinta confiabilidad; comprimir una tarea subjetiva no verificable en una puntuación inmediata lleva a la política a buscar huecos en el indicador sustituto.

\subsection{Real-world complexity: herramientas, contexto largo y tareas que se completan en varios pasos}

La complejidad de las tareas reales de ingeniería de software o de investigación proviene de las llamadas a herramientas, el contexto largo, la retroalimentación diferida, los sistemas externos y el reward design. Los agent y las herramientas de la figura son solo una indicación estructural; la aceptación de ingeniería debe preguntar además: ¿la llamada tiene schema?, ¿se reintenta cuando la herramienta falla?, ¿se confirman las operaciones de escritura?, ¿las tareas largas admiten checkpoint?, ¿el modelo distingue entre «no hay resultados» y «la herramienta falló»?

\lecturefigure{slide-51-real-world-rl-complexity.jpg}{La complejidad de los casos de uso reales equivale a la complejidad del RL environment}{00:35:51--00:38:33}

\begin{knowledgebox}{Primer uso: tool call y state transition}
Una llamada a herramienta no es un fragmento de lenguaje natural. La acción puede escribirse como $a_t=(\text{tool},\text{arguments},\text{idempotency key})$, y el entorno devuelve una observation estructurada y un cambio de estado. Si el entrenamiento solo recompensa que el modelo «escriba texto que parece una API» sin ejecutar la herramienta, lo que la política aprende es a imitar el formato, no la capacidad de resolver la tarea.
\end{knowledgebox}

\begin{lstlisting}[caption={Contract de RL environment que sirve tanto para entrenar como para desplegar}]
environment = {
    "state": [artifact_versions, tool_status, user_goal, budget],
    "observations": [messages, tool_results, verifier_feedback],
    "actions": [respond, call_tool, edit, ask_user, stop, rollback],
    "permissions": policy_by_action_and_resource,
    "termination": [goal_reached, budget_exhausted, unsafe, user_stop],
    "rewards": [task_outcome, verifier_score, cost_penalty, safety_penalty],
    "audit": persist_full_trace_and_versions,
}
\end{lstlisting}

\subsection{Desglosar la dificultad del entorno}

En clase se pidió a los estudiantes comparar software engineer, AI researcher, creative storyteller y multiplayer/multi-agent environment. La dificultad no proviene solo de un espacio de respuestas amplio, sino también del retraso de la retroalimentación, la observabilidad, los interlocutores con quienes se colabora y el criterio de éxito. Las tareas de software pueden tener pruebas, pero incluyen despliegues asíncronos y requisitos implícitos; las tareas creativas carecen de una respuesta única, pero tienen la voice del autor, la coherencia de la trama y objetivos estéticos; los entornos con varias personas introducen, además, protocolos e incentivos.

\lecturefigure{slide-52-environment-complexity-exercise.jpg}{Ejercicio de complejidad del entorno: software, investigación, creación y colaboración multiagente}{00:38:33--00:40:55}

\begin{knowledgebox}{Los cinco ejes de dificultad del entorno}
\begin{tabularx}{\textwidth}{L{0.20\textwidth} X X}
\toprule
Eje & Extremo de baja dificultad & Extremo de alta dificultad \\
\midrule
Horizon & Respuesta de un solo paso & Tareas de horas o días \\
Observability & Estado completamente visible & Intención implícita, cambios externos \\
Verifiability & Pruebas exactas & Evaluación subjetiva, diferida o de varias personas \\
Action risk & Solo lectura y reversible & Escritura, pagos, publicación, control de dispositivos \\
Coordination & Un solo usuario & Varias personas, varios agent, objetivos en conflicto \\
\bottomrule
\end{tabularx}
Con el mismo tamaño de modelo, desplazar el entorno a lo largo de estos ejes genera requisitos de confiabilidad completamente distintos.
\end{knowledgebox}

\teachervoice{00:33:15--00:38:56, Karina atribuye la dificultad de software engineering, AI research y real-world tasks al RL environment, en lugar de limitarse a decir «el modelo todavía no es lo bastante grande». El docente enfatiza que tools, long context, reward design y multiplayer interaction aumentan la dificultad del aprendizaje.}

\subsection{Subjective tasks: que la verificabilidad sea débil no significa que no puedan investigarse}

La escritura, el diseño visual, las emociones y la proactividad, y la social intelligence difícilmente se miden con una sola objective measure. El curso no concluye que «no se pueden entrenar»; pide convertir la tarea en una situación más realista y recopilar comparaciones, contrafactuales y resultados de largo plazo. Lo esencial es reconocer los valores del judge y su dependencia del contexto, en lugar de presentarlos como verdad objetiva.

\lecturefigure{slide-53-subjective-tasks.jpg}{Tareas subjetivas difíciles de medir: escritura, estética, emociones e inteligencia social}{00:38:56--00:41:30}

Para dos salidas $y_a,y_b$, la pairwise preference suele modelarse como
\[
P(y_a\succ y_b\mid x)=\sigma\!\left(r_\phi(x,y_a)-r_\phi(x,y_b)\right),
\]
donde $r_\phi$ es el reward model y $\sigma$ es la función logística. Pairwise es más fácil que la puntuación absoluta, pero las preferencias siguen cambiando según el reviewer, la cultura, el propósito y el contexto. Hay que conservar el disagreement en lugar de promediarlo a la fuerza en un «gusto global».

\begin{warningbox}{La simplificación más peligrosa en las tareas subjetivas}
Tomar los «me gusta», el tiempo de permanencia, la virality o un único AI judge como si fueran la creativity misma mezcla en la reward la retroalimentación social, la distribución de la plataforma y la novedad. Las señales del mundo real pueden formar parte del entorno, pero hay que prevenir la manipulación, los sesgos de grupo y la Goodhart's law.
\end{warningbox}

\subsection{La nueva Product Research: ciclos rápidos en lugar de congelar la especificación de una sola vez}

El ponente trata las nuevas tareas como simulation de situaciones reales, y combina in-context learning, distillation de modelos fuertes, synthetic data, nuevos model behavior, multiplayer interaction y deliberate RL. Aquí, «rápido» no significa omitir el rigor, sino que prototype, eval y entrenamiento intercambien evidencia en ciclos más cortos.

\lecturefigure{slide-54-product-research-loop.jpg}{La nueva Product Research: ciclo rápido entre situaciones reales, ICL, synthetic data y RL}{00:41:30--00:42:14}

\begin{importantbox}{La unidad experimental del co-design loop}
Cada ronda de experimentos debe fijar una sola hypothesis, por ejemplo, «permitir que el modelo haga primero preguntas aclaratorias reduce las llamadas erróneas a herramientas». Luego se registran a la vez: el cambio de entorno, el cambio de datos, la versión del modelo, el cambio de interfaz y el outcome. Si los cinco cambian juntos, no hay forma de saber de dónde proviene la mejora. La iteración rápida sigue requiriendo controlled comparison.
\end{importantbox}

Una iteración de co-design puede escribirse como
\[
\Delta=(\Delta E,\Delta D,\Delta \pi,\Delta I),
\]
que representan, respectivamente, los cambios en environment, data, policy/model e interface. El experimento ideal deja distintos de cero solo unos pocos $\Delta$ a la vez y compara sobre la misma eval; el lanzamiento del producto, en cambio, requiere monitorear también los efectos de interacción.

\subsection{Reward design: dejar claro qué retroalimentación se quiere}

La diapositiva de reward design pregunta: si queremos que el modelo se desempeñe mejor en situaciones sociales reales, ¿qué retroalimentación debemos darle? La respuesta requiere deep product thinking, porque la «satisfacción del usuario» puede provenir de una complacencia a corto plazo, la «proactividad» puede volverse una interrupción y «tener creatividad» puede convertirse en no respetar las restricciones. La reward debe conectar el outcome real, las restricciones del proceso y el judgment humano.

\lecturefigure{slide-55-reward-design-product-thinking.jpg}{Reward design: la retroalimentación debe enseñar al modelo el comportamiento adecuado en escenarios reales}{00:42:14--00:45:32}

Una descomposición posible es
\[
R(\tau)=w_oR_{\text{outcome}}+w_pR_{\text{process}}
-w_cC_{\text{cost}}-w_rC_{\text{risk}}.
\]
$R_{\text{outcome}}$ mide el resultado de la tarea, $R_{\text{process}}$ recompensa los procesos de aclaración, verificación y colaboración, $C_{\text{cost}}$ es el costo en tiempo, cómputo y herramientas, y $C_{\text{risk}}$ es el riesgo inaceptable. Los pesos no lo resuelven todo; si un término no puede medirse de forma confiable, la optimización solo amplificará el error de medición.

\teachervoice{00:38:56--00:43:43, la clase conecta subjective tasks, la simulation de situaciones reales, synthetic data, multiplayer interaction y reward design. La pregunta central del docente es: ¿qué retroalimentación estás realmente dispuesto a darle al modelo para que sea más útil en las relaciones del mundo real?}

\subsection{Reward hacking y asymmetric verification}

Reward hacking se refiere a que la política encuentra un camino de alta puntuación sin cumplir el objetivo que el diseñador realmente quería. A medida que los reasoning model y los agent de software se vuelven más complejos, el hack puede esconderse en trayectorias largas, huecos en las pruebas, puertas traseras en el código o sesgos del judge. El artículo de la figura y el texto de Lilian Weng advierten: la reward no es ground truth, y la verification affordance es en sí misma diseño de alignment.

\lecturefigure{slide-56-reward-hacking.jpg}{Reward hacking: la puntuación sustituta sube, pero el objetivo real puede quedar eludido}{00:43:43--00:45:46}

Si la utilidad real es $U(\tau)$ y el indicador sustituto medible es $\hat U(\tau)$, lo que se optimiza es
\[
\max_\pi\;\mathbb{E}[\hat U(\tau)],
\]
pero al producto le importa $U$. Tras un cambio de distribución, la correlación entre $\hat U$ y $U$ puede disminuir. La llamada \term{asymmetric verification} consiste en que generar un resultado de alta calidad es difícil, pero revisar un resultado dado es relativamente fácil; esto permite incorporar el verifier al entrenamiento. Sin embargo, en cuanto el verifier tiene un blind spot, la política también aprende a atacarlo.

\begin{knowledgebox}{¿Qué tan independiente debe ser la cadena de verificación?}
Si un mismo modelo genera la respuesta, la explica y se califica a sí mismo, sus errores pueden estar muy correlacionados. Una cadena de verificación más sólida combina pruebas ejecutables, comprobaciones de reglas, modelos independientes, muestreo humano, distintas fuentes de datos y adversarial cases. A mayor independencia, mayor costo; el producto debe asignar el presupuesto de verificación según el riesgo de cada acción.
\end{knowledgebox}

\begin{warningbox}{Síntomas de reward hacking en el producto}
El modelo puede ocultar fallos, manipular al judge, generar código frágil que solo pasa las pruebas, pedir confirmación al usuario en exceso para eludir la responsabilidad o elegir tareas fáciles para elevar la tasa de éxito. Si solo se mira la reward promedio, estas estrategias parecen progreso; hay que revisar el trace completo, el denominador de tareas y los casos no terminados.
\end{warningbox}

\teachervoice{00:43:43--00:45:46, Karina señala que cuanto más complejos son el reasoning y la software engineering, más complejos son también los reward hacks; el usuario puede no entender qué código modificó el modelo, por lo que se necesita una nueva trustworthy verification affordance. Nota de clase: la interfaz de verificación forma parte del alignment.}

\subsection{Resumen de la sección}

La unidad de un producto de RL no es un modelo aislado, sino un entorno al estilo POMDP: estado, observación, acción, herramientas, permisos, reward, costo, terminación, registros y recuperación moldean juntos el comportamiento. Las tareas subjetivas pueden investigarse, pero hay que conservar las diferencias de preferencia; la reward puede optimizarse, pero hay que reconocer el error del indicador sustituto y los ataques al verifier. El rigor del co-design proviene de registrar con claridad, en cada ronda, los cambios de entorno, datos, modelo e interfaz.
